\begin{figure}[h]
\sporkForceStacktrue
\sporkSemanticsStuff
\hrule
\begin{mathpar}
\inferrule*[right=\sporkSAllocStep]{
  \sporkRule{\sporkSub{\sporkexpr_1}{\sporkexpr_2}, \sporksigma, \sporkSub{\sporktr^*_1}{\sporktr^*_2}}
}
{
  \sporkRule{\texttt{$\sporkx$: $\texttt{barrier}$[$\sporkz^*$, $\sporkSub{\sporkexpr_1}{\sporkexpr_2}$, $\sporkexpr^*$] @ $\sporkbarMech$}, \sporksigma, \sporkSub{\sporktr^*_1}{\sporktr^*_2}}
}\\
\inferrule*[right=\sporkSAlloc]{
  \sporktr^*_2 = \sporktr^*_1 \sporkConcatChk \sporkallocTrDetail(\sporkx[\sporkz_1, \ldots, \sporkz_n], \sporkbarMech) \\
  \sporksigma' = \sporksigma[(\sporkx, \sporkz^*) \mapsto 0 \mid (\sporkx, \sporkz^*) \in \sporkWunroll(\sporkx[0\texttt{:}\sporkz_1, \ldots, 0\texttt{:}\sporkz_n])]
}
{
  \sporkRule{\sporkSub{\texttt{$\sporkx$: $\texttt{barrier}$[$\sporkz_1, \ldots, \sporkz_n$] @ $\sporkbarMech$}}{\texttt{pass}}, \sporkSub{\sporksigma}{\sporksigma'}, \sporkSub{\sporktr^*_1}{\sporktr^*_2}}
}
\\
\inferrule*[right=\sporkSArriveStep]{
  \sporkRule{\sporkSub{\sporkwH_1}{\sporkwH_2}, \sporksigma, \sporkSub{\sporktr^*_1}{\sporktr^*_2}}
}
{
  \sporkRule{\texttt{Arrive($\sporkQtlArrive$) >> $\sporkwzH^*,\sporkSub{\sporkwH_1}{\sporkwH_2},\sporkwH^*$} \#\sporktidMap, \sporksigma, \sporkSub{\sporktr^*_1}{\sporktr^*_2}}
}\\
\inferrule*[right=\sporkSAwaitStep]{
  \sporkRule{\sporkSub{\sporkwH_1}{\sporkwH_2}, \sporksigma, \sporkSub{\sporktr^*_1}{\sporktr^*_2}}
}
{
  \sporkRule{\texttt{Await($\sporkSub{\sporkwH_1}{\sporkwH_2}$, $\sporkQtlAwait$, $\sporkz_\mathrm{lag}$)} \#\sporktidMap, \sporksigma, \sporkSub{\sporktr^*_1}{\sporktr^*_2}}
}
\end{mathpar}
\caption{Synchronization Structural Semantics.
Note, $\sporksigma(\sporkx)$ is used to store ``arrive counts'' for barrier variables $\sporkx$.
This is initially 0 (\sporkSAlloc), but is overridden in some cases for mbarriers (Section~\ref{sec:SmemAppendix});
said section also explains the synchronization obligations imposed by $\sporkallocTrDetail$.
% Totally ignored: Await count assertion.
}
\label{fig:SyncStructuralSemantics}
\hrule
\end{figure}
